\documentclass[../main.tex]{subfiles}

\begin{document}


Trí tuệ nhân tạo đang thay đổi căn bản cách thức tiếp cận tri thức trong thế kỷ 21. Các mô hình ngôn ngữ lớn như GPT-4o không còn chỉ là công cụ tìm kiếm thông tin mà đã trở thành trợ lý có khả năng lý giải, phân tích và phản hồi theo từng ngữ cảnh cụ thể. Trong bối cảnh đó, ứng dụng AI vào giáo dục lập trình mở ra những khả năng mà mô hình dạy học truyền thống khó đáp ứng được.

Thực tế giảng dạy lập trình tại các trường đại học cho thấy một nghịch lý quen thuộc: lớp học có một giảng viên nhưng sinh viên thì học với tốc độ và xuất phát điểm hoàn toàn khác nhau. Người nắm vững kiến thức nền bước tiếp nhanh, trong khi người còn lỗ hổng khái niệm cơ bản dần bị bỏ lại phía sau. Ngoài giờ học chính thức, sinh viên chủ yếu tìm đến ChatGPT hay các công cụ tìm kiếm khi gặp vướng mắc. Những công cụ này cho câu trả lời tức thì nhưng hoàn toàn mù quáng trước bối cảnh của người hỏi, không biết sinh viên đang ở giai đoạn nào, đã nắm được những gì và thường xuyên vấp phải lỗi ở điểm nào.

Đề tài này xuất phát từ kỳ vọng xây dựng được một hệ thống có khả năng cải thiện những bất cập hiện nay. Thay vì hoạt động theo kiểu trả lời từng câu hỏi rời rạc, hệ thống PPAT ghi nhận toàn bộ quá trình học tập của mỗi sinh viên, từ trình độ ban đầu, các lỗi lặp lại đến tốc độ tiến bộ theo thời gian, và lấy đó làm cơ sở để điều chỉnh mọi phản hồi cho phù hợp. Thay vì đưa ra đáp án sẵn, hệ thống đặt câu hỏi gợi mở để sinh viên tự tư duy và hình thành kỹ năng giải quyết vấn đề một cách chủ động.

Về mặt kỹ thuật, hệ thống được xây dựng trên nền tảng Next.js kết hợp Monaco Editor cho phép viết và chạy code trực tiếp trên trình duyệt mà không cần cài đặt thêm công cụ. Backend sử dụng Fastify viết bằng TypeScript với PostgreSQL lưu trữ hồ sơ học tập và Redis xử lý cache. GPT-4o được tích hợp thông qua cơ chế system prompt sinh động, mỗi lần gọi API đều mang theo đầy đủ ngữ cảnh về người học. Toàn bộ hệ thống được đóng gói và vận hành qua Docker Compose.

Báo cáo được tổ chức thành 4 chương:

\begin{itemize}
    \item Chương 1 trình bày nền tảng lý thuyết về AI, LLM và các công nghệ sử dụng.
    \item Chương 2 trình bày kết quả khảo sát thực tế, xác định yêu cầu chức năng và phi chức năng, đặc tả Use Case.
    \item Chương 3 trình bày thiết kế kiến trúc SOA, cơ sở dữ liệu, AI Orchestrator và giao diện.
    \item Chương 4 trình bày quá trình triển khai, kiểm thử và tài liệu hướng dẫn sử dụng.
\end{itemize}


\end{document}